\chapter{Research Lineages and Their Afterlives}

The systems cited throughout this manual did not share one fate. Some remain
active, some became products or standards, some were absorbed into a larger
platform, and some ended after demonstrating one useful mechanism. This
appendix separates those outcomes from CharlotteOS's own inheritance.

A research operating system need not replace Unix or Windows to matter. Its
usual afterlife is selective: a fast IPC path, an authority model, a queue
discipline, or a deployment abstraction survives at a different layer while
the original system and its stronger unifying thesis disappear. Conversely,
the adoption of one mechanism does not validate every claim made by the
original architecture.

\section{Capability kernels and confined Unix}

\textbf{KeyKOS, EROS, and CapROS.} KeyKOS was a commercial capability system;
EROS carried its object-capability and persistence ideas into research, and
CapROS continues EROS as a small open-source project. They did not displace
conventional operating systems, but their account of authority---possession,
delegation, attenuation, confinement, and persistence as properties of
references rather than global names---remains foundational. The CapROS project
explicitly identifies itself as the continuation of EROS
(\url{https://www.capros.org/}).

\textbf{L4 and seL4.} This lineage achieved both assimilation and independent
survival. Earlier L4 variants demonstrated that microkernel IPC could be fast
and shipped commercially in mobile devices. seL4 added machine-checked
functional correctness and security properties, became open source in 2014,
and acquired a foundation and commercial high-assurance ecosystem in 2020.
The project's history records research prototypes, the HACMS deployment, proof
milestones, and the later ecosystem
(\url{https://sel4.systems/About/history.html}).

\textbf{Capsicum and CHERI.} Capability mechanisms also entered conventional
systems in narrower forms. Capsicum has shipped in FreeBSD since version 9.0
and treats file descriptors as attenuable capabilities inside Unix
(\url{https://man.freebsd.org/cgi/man.cgi?query=capsicum&sektion=4}). CHERI
moved bounds and permissions into tagged pointers and processor support. Arm's
Morello board and CheriBSD established a substantial experimental stack, while
CHERI-RISC-V work continues toward standardisation
(\url{https://www.cl.cam.ac.uk/research/security/ctsrd/cheri/}). These are
complementary outcomes: Capsicum confines named Unix resources and CHERI
constrains memory references; neither is by itself a complete object-capability
operating system.

\subsection{CharlotteOS inheritance}

CharlotteOS adopts the strong local authority rule: an opaque, typed object
capability is authority, may carry attenuated rights, and is not a name,
notification, completion, or unit of work. Reply tokens are linear authority;
memory, device, DMA, observer, and bootstrap access are explicitly delegated.
The present kernel does \emph{not} retain a general derivation tree, prove
confinement, or provide CHERI-style protection for arbitrary pointers. This is
an EROS/seL4-inspired object model, not an seL4 refinement or a CHERI substitute.

\section{IPC, typed channels, and ownership}

\textbf{Mach.} Pure Mach did not become the dominant Unix organisation, but
it was absorbed into Apple XNU. XNU retains Mach tasks, virtual memory, ports,
and IPC while co-locating BSD, networking, filesystems, and I/O components in
one kernel for performance. Apple's description is candid about this hybrid
outcome
(\href{https://developer.apple.com/library/archive/documentation/Porting/Conceptual/PortingUnix/additionalfeatures/additionalfeatures.html}{Apple's
XNU architecture note}).
The lesson is double-edged: message-mediated authority survived, while the
pure user-server decomposition did not.

\textbf{Singularity.} Microsoft ended the standalone research OS and evolved
its design into the internal Midori project. Midori never became a commercial
OS, although Microsoft reports that it ran part of the company's natural
language search service
(\url{https://www.microsoft.com/en-us/research/project/singularity/}). Typed
channel contracts, software-isolated processes, manifest-checked components,
and ownership transfer through an exchange heap outlived the platform as
research ideas. Managed-language enforcement did not become a mainstream OS
boundary, but ownership types and message-oriented runtimes became normal
engineering tools.

\textbf{Xous.} Xous remains a working Rust microkernel for embedded devices,
not merely an historical analogy. Its userspace servers and distinct scalar,
send, lend, and mutable-lend messages make it CharlotteOS's closest direct IPC
ancestor (\url{https://xous.dev/kernel/}).

\subsection{CharlotteOS inheritance}

CharlotteOS takes Xous's isolated-server model and its separation of scalar
from memory-bearing IPC. It takes Singularity's ownership-moving discipline
but enforces it through capabilities, page mappings, teardown state machines,
and DMA rules rather than a managed-language verifier. Copy, move, read-only
borrow, and writable borrow have explicit failure, cancellation, restoration,
and rollback behavior. Typed protocol crates provide channel contracts above
an ABI that still validates untyped words at the kernel boundary.

\section{Multicore systems, runtimes, and fast I/O}

\textbf{Barrelfish and Arrakis.} Barrelfish argued that a multicore machine
should be treated as a distributed system with explicit messages and
topology-aware state placement. The project now identifies itself as inactive
(\url{https://barrelfish.org/download.html}). Arrakis gave applications
virtualised device queues while leaving allocation and protection policy in
the kernel; it merged back into Barrelfish in 2015. The code line ended, but
per-core runtimes, hardware queue virtualisation, and control/data-plane
separation became widespread design patterns.

\textbf{IX, DPDK, and Seastar.} IX was a research dataplane OS. Its
run-to-completion, per-core, batched approach continued in datacenter runtimes.
DPDK succeeded by becoming a Linux Foundation userspace packet-processing
project with production device support (\url{https://www.dpdk.org/about/}).
Seastar remains an active shared-nothing, shard-per-core, futures runtime used
by ScyllaDB (\url{https://seastar.io/}). In this family, the ideas prospered as
libraries and queues on Linux more than as replacement operating systems.

Windows completion ports, BSD \texttt{kqueue}, Linux \texttt{epoll}, and
Linux \texttt{io\_uring} show the same selective adoption. Submission,
retained completion, aggregation, and batching became normal interfaces
without making notification, readiness, completion, IPC, and authority one
thing.

\subsection{CharlotteOS inheritance}

Sitas is the direct shard execution discipline and is Seastar-like: one owner
per shard, cooperative futures, bounded messages, and deliberate placement.
Barrelfish supplies the analogy for explicit cross-LP coordination. IX, DPDK,
and Arrakis inform userspace drivers, batched queues, IOMMU protection, and
buffer ownership. CharlotteOS makes a different authority choice from Arrakis:
devices are delegated to isolated driver services, not directly to ordinary
applications. Completion queues retain terminal operation records and remain
separate from endpoint IPC and Rust \texttt{Waker}s.

\section{Service recovery}

MINIX 3, CuriOS, and Pebble explored minimally privileged userspace services,
drivers, supervision, and restart. Nooks and SafeDrive instead tried to
contain faulty extensions while preserving compatibility with in-kernel
drivers. MINIX remains a research and teaching system; the other projects are
completed prototypes. Their broad result---isolate components and supervise
them---is now conventional. Their harder problem did not disappear: hardware,
DMA, persistent state, and externally visible effects prevent arbitrary
restart from being transparent.

CharlotteOS adopts the strong containment boundary and the honest recovery
contract. The supervisor can revoke capabilities and borrows, reset a device,
reconcile operations and DMA, start a fresh generation, and require clients to
resolve it again. A stateful upgrade transfers explicitly serialised state;
it does not inherit a process image or promise uninterrupted execution. This
is closer to MINIX's supervision goal than to transparent recovery.

\section{From distributed operating systems to cluster managers}

\textbf{Amoeba.} Amoeba's last official release was 5.3 in 1996. Its processor
pools, capability-protected objects, location-independent RPC, and FLIP network
did not become a maintained product platform
(\url{https://www.cs.vu.nl/pub/amoeba/manuals/usr.pdf}). The component ideas
did become ordinary: pooled compute, object references, RPC, directories, and
replicated services. The part modern systems mostly rejected was the promise
that a distributed computer should look like one failure-free machine.

\textbf{Plan 9 and Inferno.} Bell Labs development ended, but Plan 9 has
community descendants. Its 9P protocol remains in Linux, including a Virtio
transport (\url{https://www.kernel.org/doc/html/latest/filesystems/9p.html}).
Per-process namespaces and protocol-served resources stayed influential while
Plan 9 and Inferno themselves remained niche.

\textbf{Jini and MOSIX.} Sun's Jini became Apache River; Apache retired River
in 2022 (\url{https://attic.apache.org/projects/river.html}). The openMosix
project officially ended in 2008
(\url{https://sourceforge.net/p/openmosix/news/2008/02/openmosix-project-ends/}).
Discovery, leases, tuple spaces, automatic placement, and migration were not
lost. They reappeared in registries, brokers, serverless platforms, and
cluster schedulers. Modern systems usually reschedule a container or actor and
reconstruct service state instead of migrating an arbitrary process image.

\textbf{Borg, Omega, and Kubernetes.} Borg remains Google's internal
production cluster manager; Omega explored a concurrent control plane; and
Kubernetes directly inherited their lessons about declarative desired state,
scheduling, services, labels, and self-healing. Google's own account describes
that lineage (\url{https://kubernetes.io/blog/2015/04/borg-predecessor-to-kubernetes/}).
The cluster became the deployment target, but above commodity kernels and
containers rather than inside a distributed OS kernel.

\textbf{Focused mechanisms.} Raft succeeded as a reusable replicated-state-
machine protocol with many implementations (\url{https://raft.github.io/}).
SWIM-style gossip survived in memberlist and Consul. Nix made hash-identified
immutable software stores practical. TUF and Uptane became maintained update
security frameworks; Uptane is now a versioned standard under Linux Foundation
governance (\url{https://uptane.org/learn-more/about}). In-toto, Sigstore,
Cosign, and Notary brought supply-chain provenance and signed artifact identity
into cloud deployment workflows.

\subsection{Clean-break and compatibility counterpoints}

Solaris MC extended Unix into a multi-computer single-system image while
preserving the Solaris ABI
(\url{https://www.usenix.org/conference/usenix-1996-annual-technical-conference/solaris-mc-multi-computer-os}).
LegoOS reopened the operating-system boundary for disaggregated processor,
memory, and storage components, while deliberately retaining common Linux
system calls to ease adoption
(\url{https://www.usenix.org/conference/osdi18/presentation/shan}). These
systems make useful counterpoints to CharlotteOS. They move the operating
system across hardware boundaries while preserving a machine-compatible
application contract. CharlotteOS instead tests a smaller native contract of
signed execution objects, messages, and explicit capabilities.

The comparison sharpens the research question. Compatibility is valuable
because existing software, libraries, and operational knowledge are valuable.
A clean-break system must demonstrate that its shorter policy and execution
path compensates for the cost of adapting applications.

\section{Cluster-native placement, ingress, and lifecycle}

\subsection{Desired state and replica placement}

Borg, Omega, and Kubernetes established the modern control-plane pattern:
declare desired state, schedule from constraints, observe actual state, and
reconcile the difference. Kubernetes describes its controllers directly as
control loops
(\url{https://kubernetes.io/docs/concepts/architecture/controller/}). Ceph's
CRUSH algorithm provides a related precedent for deriving deterministic
placement from a compact topology and failure-domain map rather than storing
one decision per object.

CharlotteOS now applies these ideas to native protected domains. Its
\texttt{dns} catalog is a Raft state machine that commits membership, signed
releases, replica policy, concrete assignments, deployment generations,
operational bindings, shutdown intent, and per-node readiness. The leader
resolves singleton, fixed-replica, and every-eligible-node policy against
admitted, non-draining voters. Affinity and anti-affinity relate components,
and node agents reconcile every assigned deployment.

Placement is also an authority decision. A selected node must publish readiness
for the exact committed generation before the service can receive grants or new
cluster-addressed flows. Capacity, trust labels, named failure domains,
health-driven replacement, and rolling availability policy remain the next
scheduler work.

\subsection{Cluster service identity and Direct Server Return}

Distributed object systems established location-independent service names.
Modern software load balancers use deterministic connection mapping and
retained flow state to limit disruption when backend sets change. Maglev is a
production-scale example
(\url{https://www.usenix.org/system/files/conference/nsdi16/nsdi16-paper-eisenbud.pdf}).

CharlotteOS derives its DSR backend set directly from Raft membership,
committed application placement, exact-generation readiness, and drain intent.
The frame router consumes an immutable epoch and applies rendezvous hashing to
each five-tuple. A remote backend receives the original packet through one L2
envelope, owns the TCP connection, and replies directly. The VIP advertiser,
execution owner, and public service identity may therefore be different.

The present implementation supports a bounded IPv4/TCP configuration and an
administratively trusted cluster L2. Multiple services, hostile-fabric
authentication, larger membership, and sustained churn remain experiments.
The distinctive question is whether packet eligibility can remain a direct,
inspectable consequence of application placement and lifecycle.

\subsection{Remote calls and explicit failure}

CharlotteOS retains Amoeba's common service model and accepts the critique in
Waldo et al.'s \emph{A Note on Distributed Computing}. Relmsg v3 carries
32-bit lengths, fragmentation, request and session identity, bounded duplicate
handling, retries, and deadlines. Distributed calls include the target
generation and can report an uncertain result. They do not promise that a
network call has local latency, synchronous revocation, or a knowable outcome
after reply loss. General cryptographic bearer capabilities, transitive
distributed revocation, and durable exactly-once execution remain open.

Lifecycle uses the same generation discipline. Exact readiness gates
publication and ingress; committed drain removes a node from new work; a
signed grace period bounds cooperative service shutdown; and forced retirement
closes the domain after the deadline. Stateful cross-node movement still needs
an explicit, failure-safe transfer of mutable ownership.

\section{Operational trust and managed-service authority}

\subsection{Software supply chain and role separation}

TUF separates update roles, versions, thresholds, and expiry
(\url{https://theupdateframework.github.io/specification/v1.0.28/}). In-toto
binds artifacts to verifiable supply-chain steps
(\url{https://www.usenix.org/conference/usenixsecurity19/presentation/torres-arias}).
RFC~9180 defines Hybrid Public Key Encryption
(\url{https://www.rfc-editor.org/info/rfc9180/}). CharlotteOS composes those
ideas into a runtime-authority path:

\begin{itemize}
    \item development signs executable behavior and release descriptors;
    \item operations independently signs an HPKE-encrypted Kafka or S3 profile
        bound to one cluster, release, connector, sequence, and expiry;
    \item the leader verifies both authorities against trusted UTC and commits
        only compact references and replay fences; and
    \item the selected node retrieves the ciphertext, after which the kernel
        re-verifies it, decrypts into transient zeroizing memory, and moves a
        read-only profile capability into the connector.
\end{itemize}

This gives the two organisational roles different powers. Development cannot
select production credentials. Operations cannot replace the approved
executable through a connector binding. Production KMS/HSM custody, measured
boot or attestation, rotation, recovery, audit, and readiness-driven connector
replacement remain necessary.

\subsection{Connector-confined managed services}

SPIFFE addresses a neighbouring problem by delivering platform-selected,
short-lived identities through a privileged workload API
(\url{https://spiffe.io/docs/latest/spiffe-specs/spiffe_workload_api/}).
CharlotteOS often keeps the external identity one step farther away: Kafka and
S3 credentials remain in a connector, while an application receives an
attenuated endpoint for named routes, prefixes, or operations. Independently
named connectors allow several roles and destinations. The Kafka-step service
owns the transaction and invokes a business procedure that never receives the
Kafka credential or delivery token.

The connector can consequently act as an authority firewall and an operations-
owned integration point. Its value should be measured through credential
exposure, policy review effort, reuse, revocation latency, failure isolation,
and the effect of connector compromise.

\section{AI-assisted adaptation and the POSIX question}

Clean-break systems historically faced a severe software-adoption cost. That
explains why LegoOS retained a Linux ABI and why capability systems such as
Capsicum add confinement within Unix. Coding agents change this constraint by
reducing some of the human work required to discover assumptions, translate
APIs, generate adapters, and build tests. They do not make generated code
correct by origin.

Recent migration benchmarks show both promise and substantial failure. The
2025 CODEMENV study evaluates package and environment adaptation
(\url{https://arxiv.org/abs/2506.00894}); FreshBrew evaluates project-level
Java migration and emphasises semantic preservation under strong tests
(\url{https://arxiv.org/abs/2510.04852}). These results make compatibility cost
an empirical question, while reinforcing the need for conformance tests,
review, reproducible generation, and provenance.

Durga provides a constrained CharlotteOS experiment. A BPMN process model and
developer-owned procedure generate the Charlotte adapter, capability
requirements, resource declarations, build inputs, replica and affinity
policy, and deployment metadata. Platform mechanics remain in generated code;
the business rule remains testable on a conventional host. The research test
is whether this reduces adaptation cost without hiding authority, lifecycle,
or failure semantics.

\subsection{CharlotteOS inheritance}

CharlotteOS separates a human-readable name from the connection capability
returned by policy-controlled lookup. Its replicated catalog and relmsg path
carry request identity, generation fencing, duplicate handling, deadlines, and
explicit uncertain outcomes. Signed releases drive deterministic replica
placement; exact-generation readiness gates publication and DSR; and assigned
agents retrieve digest-pinned artifacts from central S3. Applications obtain
signed, attenuated grants, while independently authorised connector profiles
retain infrastructure details and credentials.

The result is Amoeba-like in invocation, Borg/Kubernetes-like in desired-state
control, Maglev/CRUSH-like in stable deterministic placement, and
TUF/in-toto-like in trust decomposition. CharlotteOS applies those ideas
directly to native capability-constrained execution objects. Its bounded
implementation does not yet establish production scale, general distributed
revocation, stateful migration, automatic health recovery, complete trust
lifecycle, or production administration.

\section{The synthesis and the rejected inheritances}

\begin{longtable}{|p{0.19\textwidth}|p{0.32\textwidth}|p{0.38\textwidth}|}
\hline
\textbf{Relationship} & \textbf{Source ideas} & \textbf{CharlotteOS status} \\
\hline
\endfirsthead
\hline
\textbf{Relationship} & \textbf{Source ideas} & \textbf{CharlotteOS status} \\
\hline
\endhead
Direct source & Xous IPC and memory messages; Sitas shard ownership and async
execution & Implemented locally, combined with CharlotteOS capabilities,
memory objects, completion queues, and teardown semantics \\
\hline
Adopted and changed & EROS/seL4 authority; Singularity ownership transfer;
Barrelfish messages; MINIX recovery; completion-ring I/O & Implemented in
bounded local forms; no compatibility or refinement claim \\
\hline
Focused protocol & Raft replication; rendezvous placement;
Ed25519/SHA-256 verification; RFC~9180 HPKE & Implemented and tested in bounded
forms across catalog replication, replica placement, DSR ingress, and
role-separated connector launch \\
\hline
Architectural influence & Amoeba invocation; IX/DPDK/Arrakis fast paths;
Borg/Kubernetes reconciliation; Maglev/CRUSH placement;
Nix/TUF/Uptane/in-toto trust; SPIFFE identity & Implemented cluster foundations
with production-scale policy, recovery, and trust lifecycle still open \\
\hline
Rejected or narrowed & Failure-transparent remote calls; arbitrary process
migration; language-only isolation; direct device access by ordinary apps &
Remote uncertainty is visible; migration is protocol-specific; MMU/IOMMU and
capabilities enforce isolation; drivers retain device authority \\
\hline
\end{longtable}

\endinput
